% Aqui começa o capítulo de Considerações Finais.
\chapter{Considerações Finais}\label{chp:ConsideracoesFinais}

Este trabalho teve como objetivo simular a gestão de um projeto realista, alinhado às demandas encontradas no mercado de trabalho, por meio do desenvolvimento de um sistema de gestão de casos sociais utilizando práticas ágeis e os fundamentos do PMBOK 7ª edição. Ao longo do projeto, buscou-se integrar princípios, domínios de desempenho e métodos ágeis de forma coerente, evitando abordagens meramente prescritivas ou desconectadas da realidade organizacional.

A escolha de um sistema de gestão de casos sociais mostrou-se adequada para evidenciar a complexidade inerente a projetos institucionais e de impacto social, nos quais requisitos evoluem, múltiplos stakeholders coexistem e a qualidade da informação é tão relevante quanto a entrega técnica. Esse contexto reforçou a importância de tratar o projeto como um sistema adaptativo, no qual aprendizado contínuo e feedback são elementos centrais para o sucesso.

A aplicação dos 12 princípios do PMBOK evidenciou que a gestão de projetos contemporânea vai além do controle de cronogramas e custos, enfatizando geração de valor, engajamento das partes interessadas, liderança colaborativa e responsabilidade ética. De forma complementar, a análise dos 8 domínios de desempenho permitiu estruturar a governança do projeto de maneira integrada, conectando pessoas, processos, entregas e resultados.

O uso do Scrum como abordagem ágil possibilitou lidar de forma prática com incertezas e mudanças, por meio de ciclos curtos de planejamento, execução e inspeção. O planejamento ágil, baseado em backlog priorizado, roadmap evolutivo e métricas de acompanhamento, demonstrou ser eficaz para manter previsibilidade de curto prazo sem comprometer a adaptabilidade do projeto.

A análise de riscos e incertezas reforçou a necessidade de abandonar a visão tradicional de risco como evento isolado, adotando uma postura contínua de observação e resposta. Essa abordagem mostrou-se particularmente relevante
